\documentclass[10pt]{extarticle}
\usepackage[a4paper,margin=1.5cm]{geometry}
\usepackage[T1]{fontenc}
\usepackage{amsmath,amssymb}
\usepackage{graphicx}
\usepackage[export]{adjustbox}
\usepackage{multicol}
\usepackage{enumitem}
\usepackage{booktabs,array,tabularx}
\setlength{\columnsep}{0.7cm}
\setlength{\parindent}{0pt}
\setlength{\parskip}{3pt}
\pagestyle{empty}
\makeatletter
\newcommand{\qfigdef}[4]{\@namedef{qf@file@#1}{#2}\@namedef{qf@w@#1}{#3}\@namedef{qf@wide@#1}{#4}}
\newcommand{\qf@check}[1]{\@ifundefined{qf@file@#1}{\PackageError{score2}{Unknown figure id #1}{}}{}}
\newcommand{\qf@img}[4][]{\includegraphics[width=#2\dimexpr\@nameuse{qf@w@#4}pt\relax,max width=#3,max height=0.45\textheight,#1]{\@nameuse{qf@file@#4}}}
\newif\ifqf@full
\newif\ifqf@boxed
\newsavebox{\qf@box}
\newcommand{\qf@open}{\global\qf@fulltrue\par\global\qf@boxedtrue
  \setbox\qf@box\vbox\bgroup\hsize\textwidth\linewidth\textwidth}
\newcommand{\qf@emit}[1]{\par\egroup\global\qf@boxedfalse
  \ifdim\dimexpr\ht\qf@box+\dp\qf@box\relax<#1\textheight\noindent\usebox{\qf@box}\par\else\unvbox\qf@box\fi}
\newcommand{\qfullsplit}{\ifqf@boxed\qf@emit{0.85}\fi}
\newcommand{\qf@close}{\ifqf@boxed\qf@emit{0.45}\fi\global\qf@fullfalse}
\newcommand{\qfullbegin}{\end{multicols}\qf@open}
\newcommand{\qfullend}{\qf@close\begin{multicols}{2}}
\newcommand{\qfullstart}{\qf@open}
\newcommand{\qfulldone}{\qf@close}
\newcommand{\qfig}[2][1]{\qf@check{#2}%
  \ifqf@full\qf@blk{#1}{#2}\else
  \ifnum\@nameuse{qf@wide@#2}=1 \expandafter\@firstoftwo\else\expandafter\@secondoftwo\fi
  {\end{multicols}\begin{center}\qf@img{#1}{\textwidth}{#2}\end{center}\begin{multicols}{2}}%
  {\qf@blk{#1}{#2}}\fi}
\newcommand{\qf@blk}[2]{\par\smallskip{\centering\qf@img{#1}{\linewidth}{#2}\par}\smallskip}
\newcommand{\qfigi}[2][1]{\qf@check{#2}\qf@img[valign=c]{#1}{\linewidth}{#2}}
\makeatother
\newcommand{\papersec}[1]{\par\bigskip{\centering\Large\bfseries PAPER-#1\par}\medskip}
\newcommand{\partsec}[2]{\par\medskip\noindent{\large\bfseries PART-#1: #2}\par\smallskip}
\newcommand{\sect}[1]{\par\smallskip\noindent\textbf{#1}\par}
\newcommand{\qn}[1]{\par\smallskip\noindent\textbf{#1.}~}
\newcommand{\paraq}[2]{\par\smallskip\noindent\textbf{Paragraph for Questions #1 and #2}\par}
\newenvironment{qoptions}{\begin{itemize}[leftmargin=1.8em,itemsep=1pt,topsep=2pt,label={}]}{\end{itemize}}
\newcommand{\opt}[2]{\item[(#1)] #2}
% figure registry (generated by score2_tools.py finalize)
\qfigdef{p1_phy_s1_q01_1}{SRG (Enthusiast) - Test 37 (11 Sep 2024)-img/p1_phy_s1_q01_1.png}{104.8}{0}
\qfigdef{p1_phy_s1_q04_1}{SRG (Enthusiast) - Test 37 (11 Sep 2024)-img/p1_phy_s1_q04_1.png}{108.5}{0}
\qfigdef{p1_phy_s2_q05_1}{SRG (Enthusiast) - Test 37 (11 Sep 2024)-img/p1_phy_s2_q05_1.png}{103.9}{0}
\qfigdef{p1_phy_s2_q06_1}{SRG (Enthusiast) - Test 37 (11 Sep 2024)-img/p1_phy_s2_q06_1.png}{118.5}{0}
\qfigdef{p1_phy_s2_q07_1}{SRG (Enthusiast) - Test 37 (11 Sep 2024)-img/p1_phy_s2_q07_1.png}{93.1}{0}
\qfigdef{p1_phy_s2_q08_1}{SRG (Enthusiast) - Test 37 (11 Sep 2024)-img/p1_phy_s2_q08_1.png}{215.5}{0}
\qfigdef{p1_phy_s3_q01_1}{SRG (Enthusiast) - Test 37 (11 Sep 2024)-img/p1_phy_s3_q01_1.png}{108.2}{0}
\qfigdef{p1_phy_s3_q02_1}{SRG (Enthusiast) - Test 37 (11 Sep 2024)-img/p1_phy_s3_q02_1.png}{240.5}{0}
\qfigdef{p1_phy_s3_q03_1}{SRG (Enthusiast) - Test 37 (11 Sep 2024)-img/p1_phy_s3_q03_1.png}{42.6}{0}
\qfigdef{p1_phy_s3_q04_1}{SRG (Enthusiast) - Test 37 (11 Sep 2024)-img/p1_phy_s3_q04_1.png}{108.5}{0}
\qfigdef{p1_phy_s3_q05_1}{SRG (Enthusiast) - Test 37 (11 Sep 2024)-img/p1_phy_s3_q05_1.png}{129.3}{0}
\qfigdef{p1_phy_s3_q06_1}{SRG (Enthusiast) - Test 37 (11 Sep 2024)-img/p1_phy_s3_q06_1.png}{149.1}{0}
\qfigdef{p1_phy_s3_q07_1}{SRG (Enthusiast) - Test 37 (11 Sep 2024)-img/p1_phy_s3_q07_1.png}{82.0}{0}
\qfigdef{p1_phy_s3_q08_1}{SRG (Enthusiast) - Test 37 (11 Sep 2024)-img/p1_phy_s3_q08_1.png}{108.5}{0}
\qfigdef{p1_chem_s1_q02_1}{SRG (Enthusiast) - Test 37 (11 Sep 2024)-img/p1_chem_s1_q02_1.png}{108.5}{0}
\qfigdef{p1_chem_s2_q05_1}{SRG (Enthusiast) - Test 37 (11 Sep 2024)-img/p1_chem_s2_q05_1.png}{71.5}{0}
\qfigdef{p1_chem_s3_q02_1}{SRG (Enthusiast) - Test 37 (11 Sep 2024)-img/p1_chem_s3_q02_1.png}{88.2}{0}
\qfigdef{p1_chem_s3_q04_1}{SRG (Enthusiast) - Test 37 (11 Sep 2024)-img/p1_chem_s3_q04_1.png}{181.0}{0}

\begin{document}
\begin{center}{\Large\bfseries SRG (Enthusiast) --- Test 37 (11 Sep 2024)}\end{center}
\begin{multicols}{2}
\partsec{1}{PHYSICS}
\sect{Section-I(i)}
\qn{1} In the shown AC circuit phase difference between currents $I_1$ and $I_2$ is :
\qfig{p1_phy_s1_q01_1}
\begin{qoptions}
\opt{A}{$\dfrac{\pi}{2}-\tan^{-1}\dfrac{X_L}{R}$}
\opt{B}{$\tan^{-1}\dfrac{X_L+X_C}{R}$}
\opt{C}{$\dfrac{\pi}{2}+\tan^{-1}\dfrac{X_L}{R}$}
\opt{D}{$\tan^{-1}\dfrac{X_L-X_C}{R}+\dfrac{\pi}{2}$}
\end{qoptions}
\qn{2} In a series LCR circuit, the inductive reactance ($X_L$) is 10\,$\Omega$ and the capacitive reactance ($X_C$) is 4\,$\Omega$. The resistance ($R$) in the circuit is 6\,$\Omega$. The power factor of the circuit is :
\begin{qoptions}
\opt{A}{$\dfrac12$}
\opt{B}{$\dfrac{1}{2\sqrt2}$}
\opt{C}{$\dfrac{1}{\sqrt2}$}
\opt{D}{$\dfrac{\sqrt3}{2}$}
\end{qoptions}
\qn{3} In a series L--R growth circuit, if maximum current and maximum induced emf in an inductor of inductance 3\,mH are 2\,A and 6\,V respectively, then the time constant of the circuit is :
\begin{qoptions}
\opt{A}{1 m sec.}
\opt{B}{1/3 m/sec.}
\opt{C}{1/6 m/sec}
\opt{D}{1/2 m/sec}
\end{qoptions}
\qn{4} There exist uniform time varying magnetic field in cylindrical region of radius $R$ as shown in figure.
\qfig{p1_phy_s1_q04_1}
\textbf{Statement--1 :} Work done by induced electric field in moving a point charge from A to B is different for two paths shown.\\
\textbf{Statement--2 :} Induced electric field is a non conservative electric field.
\begin{qoptions}
\opt{A}{Statement-1 is True, Statement-2 is True; Statement-2 is a correct explanation for Statement-1}
\opt{B}{Statement-1 is True, Statement-2 is True; Statement-2 is NOT a correct explanation for Statement-1}
\opt{C}{Statement-1 is True, Statement-2 is False}
\opt{D}{Statement-1 is False, Statement-2 is True.}
\end{qoptions}
\sect{Section-I(ii)}
\qn{5} Consider a network shown in the figure consisting of a resistance, a capacitor, an inductor and three alternating voltage sources 1, 2 and 3. Terminal voltages of the sources 1, 2 and 3 are $v_1=V_0\sin(\omega t)$, $v_2=V_0\sin(\omega t+120^\circ)$ and $v_3=V_0\sin(\omega t+240^\circ)$ respectively and moduli of the reactance of the capacitor and the inductor are equal to the resistance. Then
\qfig{p1_phy_s2_q05_1}
\begin{qoptions}
\opt{A}{Peak value of voltage at point P is $v_0\left(\sqrt3+1\right)$}
\opt{B}{Peak value of voltage at point P is $v_0\left(\sqrt3-1\right)$}
\opt{C}{$v_P=v_0\left(\sqrt3+1\right)\sin(\omega t+90)$}
\opt{D}{$v_P=v_0\left(\sqrt3-1\right)\sin(\omega t+180)$}
\end{qoptions}
\qn{6} At time $t=0$, terminal A in the circuit shown in the figure is connected to B by a key and an alternating current $I(t)=I_0\cos(\omega t)$, with $I_0=1$\,A and $\omega=500\,\mathrm{rad\,s^{-1}}$ starts flowing in it with the initial direction shown in the figure. At $t=\dfrac{7\pi}{6\omega}$, the key is switched from B to D. Now onwards only A and D are connected. A total charge $Q$ flows from the battery to charge the capacitor fully. If $C=20\,\mu$F, $R=10\,\Omega$ and the battery is ideal with emf of 50\,V, identify the correct statement(s)
\qfig{p1_phy_s2_q06_1}
\begin{qoptions}
\opt{A}{Magnitude of the maximum charge on the capacitor before $t=\dfrac{7\pi}{6\omega}$ is $1\times10^{-3}$\,C.}
\opt{B}{The current in the left part of the circuit just before $t=\dfrac{7\pi}{6\omega}$ is clockwise.}
\opt{C}{Immediately after A is connected to D, the current in $R$ is 10\,A.}
\opt{D}{$Q=2\times10^{-3}$\,C}
\end{qoptions}
\qn{7} The figure shows a graph between velocity and displacement (from mean position) of a particle performing SHM
\qfig{p1_phy_s2_q07_1}
\begin{qoptions}
\opt{A}{The time period of the particle is 1.57\,s}
\opt{B}{The maximum acceleration will be 40\,cm/s$^2$}
\opt{C}{The velocity of particle is $2\sqrt{21}$\,cm/s when it is at a distance 1\,cm from the mean position.}
\opt{D}{Minimum acceleration is 4\,cm/s$^2$}
\end{qoptions}
\qn{8} A particle is executing SHM between points $-X_m$ and $X_m$, as shown in figure-I. The velocity $V(t)$ of the particle is partially graphed and shown in figure-II. Two points A and B corresponding to time $t_1$ and time $t_2$ respectively are marked on the $V(t)$ curve.
\qfig{p1_phy_s2_q08_1}
\begin{qoptions}
\opt{A}{At time $t_1$, it is going towards $X_m$.}
\opt{B}{At time $t_1$, its speed is decreasing.}
\opt{C}{At time $t_2$, its position lies in between $-X_m$ and O.}
\opt{D}{The phase difference $\Delta\phi$ between points A and B must be expressed as $90^\circ<\Delta\phi<180^\circ$.}
\end{qoptions}
\qn{9} A particle of mass $m$ is attached to one end of a mass-less spring of force constant $k$, lying on a frictionless horizontal plane. The other end of the spring is fixed. The particle starts moving horizontally from its equilibrium position at time $t=0$ with an initial velocity $u_0$. When the speed of the particle is $0.5\,u_0$, it collides elastically with a rigid wall. After this collision :-
\begin{qoptions}
\opt{A}{The speed of the particle when it returns to its equilibrium position is $u_0$}
\opt{B}{The time at which the particle passes through the equilibrium position for the first time is $t=\pi\sqrt{\dfrac mk}$}
\opt{C}{The time at which the maximum compression of the spring occurs is $t=\dfrac{4\pi}{3}\sqrt{\dfrac mk}$}
\opt{D}{The time at which the particle passes through the equilibrium position for the second time is $t=\dfrac{5\pi}{3}\sqrt{\dfrac mk}$}
\end{qoptions}
\qn{10} A particle of mass $m$ is attached to one end of a mass-less spring of force constant $k$, lying on a frictionless horizontal plane. The other end of the spring is fixed. The particle starts moving horizontally from its equilibrium position at time $t=0$ with an initial velocity $u_0$. When the speed of the particle is $0.5\,u_0$. It collides elastically with a rigid wall. After this collision,
\begin{qoptions}
\opt{A}{The speed of the particle when it returns to its equilibrium position is $u_0$.}
\opt{B}{The time at which the particle passes through the equilibrium position for the first time is $t=\pi\sqrt{\dfrac mk}$}
\opt{C}{The time at which the maximum compression of the spring occurs is $t=\dfrac{4\pi}{3}\sqrt{\dfrac mk}$}
\opt{D}{The time at which the particle passes through the equilibrium position for the second time is $t=\dfrac{5\pi}{3}\sqrt{\dfrac mk}$}
\end{qoptions}
\sect{Section-II}
\qn{1} A rectangular tank having base $15\,\mathrm{cm}\times20\,\mathrm{cm}$ is filled with water (density $\rho=1000\,\mathrm{kg\,m^{-3}}$) upto 20\,cm height. One end of an ideal spring of natural length $h_0=20$\,cm and force constant $K=280\,\mathrm{N\,m^{-1}}$ is fixed to the bottom of a tank so that spring remains vertical. This system is in an elevator moving downwards with acceleration $a=2\,\mathrm{m\,s^{-2}}$. A cubical block of side $\ell=10$\,cm and mass $m=2$\,kg is gently placed over the spring and released gradually, as shown in figure. If block is slightly pushed down from equilibrium position and released, and the frequency of its vertical oscillations is $\dfrac{a\sqrt2}{\pi}$, then calculate $a$.
\qfig{p1_phy_s3_q01_1}
\qn{2} If this body rolls purely then time period of its small oscillations is $T$. Its moment of inertia about CM is 9\,kg.m$^2$. Fill value of $\dfrac{40}{\pi}T$\,?
\qfig{p1_phy_s3_q02_1}
\qn{3} A metre stick swinging in vertical plane about an fixed horizontal axis passing through its one end undergoes small oscillation of frequency $2\sqrt2$\,Hz. If the bottom half of the stick were cut off, then its new frequency (in Hz) of small oscillation would become:
\qfig{p1_phy_s3_q03_1}
\qn{4} When the system shown in the diagram is in equilibrium, the right spring is stretched by 1\,cm. The coefficient of static friction between the blocks is 0.3. There is no friction between the bottom block and the supporting surface. The force constants of the springs are 150\,N/m and 450\,N/m (refer figure). The blocks have equal mass of 2\,kg each. Find the maximum amplitude (in cm) of the oscillations of the system shown in the figure that does not allow the top block to slide on the bottom.
\qfig{p1_phy_s3_q04_1}
\qn{5} A uniform plank of mass $m=1$\,kg free to move in the horizontal direction only, is placed at the top of a solid cylinder of mass $2m$ and radius $R$. The plank is attached to a fixed wall by means of a light spring of spring constant $k=7$\,N/m. There is no slipping between the cylinder and the plank, and between the cylinder and the ground. The time period of small oscillations of the system in sec (to the nearest integer) is
\qfig{p1_phy_s3_q05_1}
\qn{6} In the circuit shown current is changing at the rate of $\dfrac{di}{dt}=2$\,A/s. If $L_1=L_2=1$\,H and $M_{12}=\dfrac12$\,H, find emf $V$ (in volt) between terminals X \& Y.
\qfig{p1_phy_s3_q06_1}
\qn{7} The A.C. circuit is shown in figure. The frequency of the AC voltage source so that current through the source will be in same phase as of voltage of source is $x\sqrt{\dfrac{1}{LC}-\dfrac{R^2}{L^2}}$. The value of $2\pi x$ is
\qfig{p1_phy_s3_q07_1}
\qn{8} Two capacitors $C_1=C_2=\dfrac{1}{\pi^2}\times10^{-2}$\,F and inductor $L=2\times10^{-2}$\,H are connected in series as shown in the figure. Initially charge on each capacitors are $4\sqrt3\,\mu$C. At $t=0$ switch $S_1$ is closed and $t=1/400$ sec, switch $S_2$ is also closed. The maximum charge on capacitor $C_2$ during LC oscillation is $k\sqrt2\,\mu$C. Find $k$.
\qfig{p1_phy_s3_q08_1}
\partsec{2}{CHEMISTRY}
\sect{Section-I(i)}
\qn{1} For a reaction scheme $\mathrm{A}\xrightarrow{k_1}\mathrm{B}\xrightarrow{k_2}\mathrm{C}$, if the rate of formation of B is set to be zero then the concentration of B is given by :
\begin{qoptions}
\opt{A}{$\left(\dfrac{k_1}{k_2}\right)[\mathrm A]$}
\opt{B}{$(k_1+k_2)[\mathrm A]$}
\opt{C}{$k_1k_2[\mathrm A]$}
\opt{D}{$(k_1-k_2)[\mathrm A]$}
\end{qoptions}
\qn{2} Match the graphical study with the order of the reactions :
\qfig{p1_chem_s1_q02_1}
P : First order\\
Q : Second order\\
R : Zero order\\
$(a-x)$ = concentration of a reactant at time $t$.
\begin{qoptions}
\opt{A}{I--P, II--Q, III--R}
\opt{B}{I--Q, II--R, III--P}
\opt{C}{I--R, II--Q, III--P}
\opt{D}{I--R, II--P, III--Q}
\end{qoptions}
\qn{3} A gaseous mixture containing $\mathrm{N_2}$ and an unknown gas $\mathrm{A_2(g)}$ are present in a sealed vessel at 700 mm of Hg. Gas $\mathrm{A_2}$ decomposes as $\mathrm{A_2(g)\to2A(g)}$ with progress of time and decomposition follow first order kinetics. If the total pressure inside the flask measured after 30 min and 60 min are 760 and 800 mm of Hg, calculate total pressure after 80 min. [$\log2=0.30$, $\log3=0.48$, Antilog$(0.479)=3$]
\begin{qoptions}
\opt{A}{820 mm Hg}
\opt{B}{830 mm Hg}
\opt{C}{840 mm Hg}
\opt{D}{850 mm Hg}
\end{qoptions}
\qn{4} For gaseous elementary reaction ; $2\mathrm{A(g)}+\mathrm{B(g)}\to\mathrm{P(g)}$\\
Initially A and B are present in stoichiometric mole ratio and total pressure decrease from 120 torr to 80 torr in 5 minutes, than average rate of reaction is .
\begin{qoptions}
\opt{A}{8 torr/min.}
\opt{B}{6 torr/min.}
\opt{C}{2 torr/min.}
\opt{D}{4 torr/min.}
\end{qoptions}
\sect{Section-I(ii)}
\qn{5} For parallel reaction (all following first order kinetics)
\qfig{p1_chem_s2_q05_1}
at time $t=0$, initial mole of A is 1. If overall half life of reaction is 15 days and $k_1:k_2:k_3=4:2:1$, then [$\log2=0.3$, $2^{1/3}=5/4$]
\begin{qoptions}
\opt{A}{Number of moles of C after 45 days are 2.}
\opt{B}{Number of moles of B after 30 days are $\dfrac{36}{28}$.}
\opt{C}{Total moles at the end of reaction are 3.}
\opt{D}{Ratio of moles of B \& D at any time (except $t=0$) is 4.}
\end{qoptions}
\qn{6} For a first order reaction carried out at constant P \& T : $\mathrm{A(aq.)\to2B(g)+3C(g)+E(\ell)}$.\\
If volume of gas B collected at $t=1.386$ min is 200\,m$\ell$ and volume of gas C collected after sufficient time is 600\,m$\ell$ then
\begin{qoptions}
\opt{A}{Rate constant of reaction is $0.5\ \mathrm{min^{-1}}$}
\opt{B}{Volume of gas C collected after 2.772 min is 450 ml}
\opt{C}{Ratio of volume of gas C collected at $t=1.386$ min and volume of gas B collected after sufficient time is 3/4.}
\opt{D}{Half life time of reaction is 1.386 min.}
\end{qoptions}
\qn{7} $2\mathrm{X(g)}+\mathrm{Y(g)}+3\mathrm{Z(g)}\to$ Products\\
The rate equation of above reaction is given by\\
Rate $=K[\mathrm X]^1[\mathrm Y]^0[\mathrm Z]^2$\\
Choose the correct statement(s)
\begin{qoptions}
\opt{A}{If 75\% of X undergoes reaction in 20 sec, 50\% of X will react in 10 sec if $[\mathrm Z]\gg[\mathrm X]$}
\opt{B}{Rate of reaction decreases by reducing the concentration of Y to half of the original value}
\opt{C}{The half life of Z increases by increasing its concentration if $[\mathrm X]\gg[\mathrm Z]$}
\opt{D}{On doubling the concentration of X, Y \& Z, rate of reaction becomes 8 times}
\end{qoptions}
\qn{8} Radionucleide `P' decays into `Q' and `Q' decays into `R' as :\\
$\mathrm{P}\xrightarrow{t_{1/2}=1000\,\mathrm{hrs}}\mathrm{Q}\xrightarrow{t_{1/2}=10\,\mathrm{min}}\mathrm{R}$\\
Number of nuclei of Q and R after 1000 hr, if we start with 2 mol P, are. [$N_A=6\times10^{23}$]
\begin{qoptions}
\opt{A}{$N_Q=6\times10^{23}$}
\opt{B}{$N_Q=10^{20}$}
\opt{C}{$N_R=6\times10^{23}$}
\opt{D}{$N_R=10^{23}$}
\end{qoptions}
\qn{9} The correct statement(s) regarding boron nitride (BN) formation is/are :
\begin{qoptions}
\opt{A}{Boron on exposure to $\mathrm{NH_3}$ at appropriate temperature and pressure produces boron nitride}
\opt{B}{$\mathrm{B_2H_6}$ on exposure to $\mathrm{NH_3}$ at temperature less than $200^\circ$C produces boron nitride}
\opt{C}{Liquid $\mathrm{B_2O_3}$ on exposure to $\mathrm{NH_3}$ at high temperature produces boron nitride}
\opt{D}{Fluoroboric acid on exposure to $\mathrm{NH_3}$ produces boron nitride}
\end{qoptions}
\qn{10} For the given reaction select correct statement :-\\
I. $\mathrm{H_3BO_3+NaOH_{aq.}}\longrightarrow$\\
II. $\mathrm{B_2H_6+C_2H_5OH}\longrightarrow$\\
III. $\mathrm{H_3BO_{3\,aq.}}$ + catecol $\longrightarrow$\\
IV. $\mathrm{BF_3+KF}\longrightarrow$\\
V. $\mathrm{NH_4Cl+Na[BH_4]}\xrightarrow{\Delta}$ final product with respect to reactant
\begin{qoptions}
\opt{A}{Hybridisation of boron will change in I, III, IV \& V}
\opt{B}{Hybridisation of boron will change in II, IV, \& V}
\opt{C}{One of the product in reaction V is an aromatic compound}
\opt{D}{In reaction I boron containing compound in reactant act as lewis acid}
\end{qoptions}
\sect{Section-II}
\qn{1} Find the number of correct statements about $\mathrm{B_2H_6}$.\\
(i) It is lewis acidic gas.\\
(ii) It gives $\mathrm{H_2}$ gas with $\mathrm{C_2H_5OH}$.\\
(iii) It gives borazine with $\mathrm{NH_3}$ on heating (not very strong heating).\\
(iv) It hydrolysed easily with $\mathrm{H_2O}$ and gives boric acid.\\
(v) All the adjacent bond angles are identical to each other.
\qn{2} Calculate value of $|x+y-z|$ for the following silicate $[\mathrm{Si_xO_{y+z}}]^{z-}$ anion.
\qfig{p1_chem_s3_q02_1}
\qn{3} Three moles of $\mathrm{B_2H_6}$ are completely reacted with methanol. The number of moles of boron containing product formed is
\qn{4} For the reaction :
\qfig{p1_chem_s3_q04_1}
If initially 1 mol of A is taken \& finally ratio of B \& D is 6 : 1 what is total enthalpy change (cal) per mole of A reacted.
\begin{minipage}{\linewidth}\qn{5} Consider the kinetic data given in the following table for the reaction $\mathrm{A+B+C}\longrightarrow$ Product.
\begin{center}\footnotesize\setlength{\tabcolsep}{3pt}
\begin{tabular}{|p{0.14\linewidth}|p{0.14\linewidth}|p{0.14\linewidth}|p{0.14\linewidth}|p{0.2\linewidth}|}
\hline
Exp. No. & $[\mathrm A]$ (mol dm$^{-3}$) & $[\mathrm B]$ (mol dm$^{-3}$) & $[\mathrm C]$ (mol dm$^{-3}$) & Rate of reaction (mol dm$^{-3}$ s$^{-1}$)\\
\hline
1 & 0.2 & 0.1 & 0.1 & $6.0\times10^{-5}$\\
2 & 0.2 & 0.2 & 0.1 & $6.0\times10^{-5}$\\
3 & 0.2 & 0.1 & 0.2 & $1.2\times10^{-4}$\\
4 & 0.3 & 0.1 & 0.1 & $9.0\times10^{-5}$\\
\hline
\end{tabular}
\end{center}
The rate of the reaction for $[\mathrm A]=0.15$ mol dm$^{-3}$, $[\mathrm B]=0.25$ mol dm$^{-3}$ and $[\mathrm C]=0.15$ mol dm$^{-3}$ is found to be $Y\times10^{-5}$ mol dm$^{-3}$ s$^{-1}$. The value of $Y$ is \underline{\hspace{1.5cm}}.
\end{minipage}
\qn{6} $4\mathrm{A(g)}\to3\mathrm{B(g)}+2\mathrm{C(g)}$\\
At $t=0$ ; $P_A=4000$ mm of Hg; $P_B=P_C=0$\\
Rate constant of reaction $=5$ mm/sec\\
Find the total final pressure of gas mixture at $t=250$ second (in mm Hg)
\qn{7} A chemical reaction occur at the same rate at 500 K as at 300 K in presence of catalyst. If catalyst decreases the activation energy by 35 kJ/mole then calculate the activation energy of uncatalysed reaction (in kJ/mole)
\qn{8} A given sample of milk turns sour at room temperature ($27^\circ$C) in 16 hours. In a refrigerator at $7^\circ$C, milk can be stored three times as long before it sours. How long (in hours) it takes milk to sour at $57^\circ$C ?\\{[}Given : $\log3=0.5$, $10^{(7/11)}=10^{0.64}=4$]
\partsec{3}{MATHEMATICS}
\sect{Section-I(i)}
\qn{1} Let $f:\mathbb R\to[4,\infty)$ be an onto quadratic function whose leading coefficient is 1, such that $f'(x)+f'(2-x)=0$, then the value of $\displaystyle\int_1^3\dfrac{dx}{f(x)}$ is equal to
\begin{qoptions}
\opt{A}{$\dfrac14$}
\opt{B}{$\dfrac18$}
\opt{C}{$\dfrac\pi4$}
\opt{D}{$\dfrac\pi8$}
\end{qoptions}
\qn{2} For $x>-1$, let $f(x)=\displaystyle\int_0^{\pi/4}\log_e(1+x\tan z)\,dz$. Then value of $f\left(\dfrac12\right)-f\left(\dfrac13\right)$ equals
\begin{qoptions}
\opt{A}{$\dfrac\pi4\log_e\left(\dfrac98\right)$}
\opt{B}{$\dfrac\pi8\log_e\left(\dfrac98\right)$}
\opt{C}{$\dfrac\pi9\log_e\left(\dfrac89\right)$}
\opt{D}{$\dfrac\pi8\log_e\left(\dfrac32\right)$}
\end{qoptions}
\qn{3} If all roots of the equation $f(x)=x^6-12x^5+bx^4+cx^3+dx^2+ex+64=0$ are positive, then which of the following has greatest numerical (absolute) value ?
\begin{qoptions}
\opt{A}{$b$}
\opt{B}{$c$}
\opt{C}{$d$}
\opt{D}{$e$}
\end{qoptions}
\qn{4} Let $f(x)=x^2+px+q$ and $g(x)=x^2+rx+s$ where real roots of $f(x)=0$ are $\alpha,\beta$ and real roots of $g(x)=0$ are $\alpha+\delta,\beta+\delta$ such that minimum value of $f(x)$ be $-\dfrac14$ and least value of $g(x)$ occurs at $x=\dfrac72$, then correct option is $-$
\begin{qoptions}
\opt{A}{Least value of $g(x)$ is $\dfrac12$}
\opt{B}{Least value of $g(x)$ is $-\dfrac12$}
\opt{C}{Roots of the equation $g(x)=0$ are 3, 4}
\opt{D}{Roots of the equation $g(x)=0$ are $\dfrac32,\ \dfrac72$}
\end{qoptions}
\sect{Section-I(ii)}
\qn{5} Let $F(x)$ be a six degree polynomial with leading coefficient $k$. If $y=F(x)$ touches $f(x)=x+\alpha$ at $x=1,2$ and 3, then which of the following is/are CORRECT?
\begin{qoptions}
\opt{A}{If $\displaystyle\int_1^3(F(x)-f(x))\,dx=\dfrac{32}{105}$, then $k=2$}
\opt{B}{$F(x)-f(x)$ is divisible by $(x-1)^2(x-2)^2$}
\opt{C}{If $\displaystyle\int_1^3(F(x)-f(x))\,dx=\dfrac{32}{105}$, then $k=1$}
\opt{D}{$F(x)-f(x)$ is divisible by $(x-2)(x-3)^2$}
\end{qoptions}
\qn{6} Which is/are correct?
\begin{qoptions}
\opt{A}{$\displaystyle\int_0^\pi\dfrac{\sin^3x\,dx}{1+x}=3\int_0^\pi\sin2x\,(x+1)\left\{\ln\left(\dfrac{x+1}{e}\right)\right\}dx$}
\opt{B}{$\displaystyle\int_0^\pi\dfrac{\sin^3x\,dx}{2x-\pi}=\int_0^{\pi/2}\dfrac{\cos^3x}{x}\,dx$}
\opt{C}{$\displaystyle\int_0^\pi\dfrac{\sin^3x\,dx}{1+x}+\int_0^\pi\dfrac{\cos^3x}{x}\,dx=0$}
\opt{D}{$\displaystyle\int_0^\pi\dfrac{\sin^3x\,dx}{1+x}=3\int_0^\pi(\sin2x)(x+1)\ln(x+1)\,dx$}
\end{qoptions}
\qn{7} Let $I_n=\displaystyle\int_0^\pi\dfrac{\sin^2(nx)}{\sin^2x}\,dx$, $n\in\mathbb N$, then
\begin{qoptions}
\opt{A}{$I_{n+2}-2I_{n+1}+I_n=0$}
\opt{B}{$I_n=I_{n+1}$}
\opt{C}{$I_n=n\pi$}
\opt{D}{$I_1,I_2,I_3,\ldots,I_n$ are in Harmonic progression}
\end{qoptions}
\qn{8} If $\alpha\ (\alpha>0)$ is a root of $x^{11}-x^{10}+x^8-x^7+x^5-x^4+x^2-x-20=0$ then $\alpha^{12}$ is greater than
\begin{qoptions}
\opt{A}{59}
\opt{B}{61}
\opt{C}{58}
\opt{D}{57}
\end{qoptions}
\qn{9} If roots of the equation $2x^3+4x^2-5x+8=0$ are $\alpha,\beta,\gamma$ then
\begin{qoptions}
\opt{A}{$\sqrt{(2-\alpha)(2-\beta)(\gamma-2)(1-\alpha^2)(1-\beta^2)(1-\gamma^2)}=22.5$}
\opt{B}{$1+\alpha,\,1+\beta,\,1+\gamma$ are roots of $2x^3-2x^2-7x+15=0$}
\opt{C}{$\sqrt{(2-\alpha)(2-\beta)(2-\gamma)(1-\alpha^2)(1-\beta^2)(1-\gamma^2)}=22.5$}
\opt{D}{$1+\alpha,\,1+\beta,\,1+\gamma$ are roots of $2x^3-6x^2-7x-15=0$}
\end{qoptions}
\qn{10} Let $a$, $b$ and $c$ be the three distinct non-zero numbers satisfying the system of equation $\dfrac1a+\dfrac1{a-1}+\dfrac1{a-2}=1,\ \dfrac1b+\dfrac1{b-1}+\dfrac1{b-2}=1,\ \dfrac1c+\dfrac1{c-1}+\dfrac1{c-2}=1$, then
\begin{qoptions}
\opt{A}{$a+b+c=6$}
\opt{B}{$abc=2$}
\opt{C}{$(1-a)(1-b)(1-c)=1$}
\opt{D}{$(a-2)(b-2)(c-2)=2$}
\end{qoptions}
\sect{Section-II}
\qn{1} Given $\displaystyle\sum_{n=1}^{\infty}\dfrac1{n^2}=\dfrac{\pi^2}{6}$ and $\displaystyle\int_0^1\dfrac{\log_ex}{1-x^2}\,dx=-\dfrac{\pi^2}{\lambda}$, then the value of $\lambda$ equals
\qn{2} Let $\displaystyle\int_0^1\dfrac{1}{\sqrt{1+x^6}}\,dx=a$ then the value of $[e^a]$ where $[.]$ represents greatest integer function, is
\qn{3} Evaluate : $\displaystyle\lim_{x\to+\infty}\dfrac{d}{dx}\int_{2\sin\frac1x}^{3\sqrt x}\dfrac{3t^4+1}{(t-3)(t^2+3)}\,dt$
\qn{4} If $I=\displaystyle\int_1^2\dfrac{x^3-3x^2}{\ln(x-1)}\,dx$ and $J=\displaystyle\int_2^3\dfrac{(3x-5)}{\ln(x-2)}\,dx$, then $e^{I+J}$ equals............
\qn{5} If the $\displaystyle\int_0^{\pi/2}\sin\theta\,\log\sin\theta\,d\theta=\log(A/e)$ then $A$ is equal to
\qn{6} If $x^2-kx+k^2-3k=0$ \& $x^3+x^2-4x+a=0$ have two roots in common, then $a^2+k^2$ is equal to
\qn{7} Let $\alpha,\beta$ $(\alpha>\beta)$ be the roots of the quadratic equation $x^2-x-4=0$. If $P_n=\alpha^n-\beta^n$, then $\dfrac{P_{15}P_{16}-P_{14}P_{16}-P_{15}^2+P_{14}P_{15}}{P_{13}P_{14}}$ is equal to \underline{\hspace{1.2cm}}.
\qn{8} If $\alpha$ and $\beta$ are the roots of $x^4+x^3-1=0$, then $\alpha\beta$ is root of $x^6+ax^4+bx^3+cx^2+d=0$, the value of $a+b-c-d$ is\underline{\hspace{1.2cm}}.
\end{multicols}
\end{document}
